% --- Archon physics formalization source begin ---
% archon:physics
% archon:covers PhyXMiniProblems/problem_phyx_mini_0834.lean
% archon:source-report reports/phyx_mini/problem_phyx_mini_0834.source.json

\chapter{Physics problem phyx\_mini\_0834}
\label{ch:PhyXMiniProblems_problem_phyx_mini_0834}

\paragraph{Problem source.}
\#\# Question

What is the magnitude of the strength of the electric field at the position indicated by the dot?

\#\# Answer choices

- A: A: 1.3 \textbackslash{}times 10\textasciicircum{}\{5\}N/C

- B: B: 5.3 \textbackslash{}times 10\textasciicircum{}\{4\}N/C

- C: C: 1.0 \textbackslash{}times 10\textasciicircum{}\{5\}N/C

- D: D: 8.6 \textbackslash{}times 10\textasciicircum{}\{4\}N/C

\#\# Figure caption (auxiliary; use the image as primary evidence)

The image depicts a rectangular setup with three electrical charges and one point:

1. **Top Left**: A light teal circle with a "-" symbol and labeled “-10 nC.”
2. **Top Right**: A light red circle with a "+" symbol and labeled “10 nC.”
3. **Bottom Right**: A light teal circle with a "-" symbol and labeled “-5.0 nC.”
4. **Bottom Left**: A black dot representing a point with no charge label.

**Text and Measurements:**
- The top side of the rectangle is labeled as “5.0 cm.”
- The left side of the rectangle is labeled as “3.0 cm.”

The charges and points are connected with dashed lines indicating the sides of the rectangle. The colors and symbols indicate the charge polarity: negative for teal circles and positive for the red circle.

\#\# Dataset metadata

Electrostatics; Physical Model Grounding Reasoning, Spatial Relation Reasoning

\paragraph{Recorded answer/context.}
D: D: 8.6 \textbackslash{}times 10\textasciicircum{}\{4\}N/C

\paragraph{Figure/image path.}
/root/proposal\_for\_physic/science-mango/phyx\_mini\_run/phyx\_data/test\_image/834.png

\paragraph{Formalization target.}
create a compiling Lean file with sorry bodies at `PhyXMiniProblems/problem\_phyx\_mini\_0834.lean`. The Lean declarations must preserve the physical quantities, dimensions or dimensional roles, figure labels, governing-law hypotheses, and final relation expressed by this problem.
Use Mathlib/Physlib names found through LeanExplore where available. If a physics API is missing, introduce faithful local abstractions rather than scalar placeholder aliases.

\begin{theorem}[Physics formalization target]
\label{thm:physics:phyx_mini_0834:target}
\lean{PhyXMiniProblems.ProblemPhyXMini0834.problem_phyx_mini_0834}
\uses{def:physics:phyx-mini-0834:phyxminiproblems-problemphyxmini0834-fieldmagnitudeinnewtonspercoulomb, def:physics:phyx-mini-0834:phyxminiproblems-problemphyxmini0834-rectangularpointchargesetup, def:physics:phyx-mini-0834:phyxminiproblems-problemphyxmini0834-matchessuppliedchargerectanglefigure, def:physics:phyx-mini-0834:phyxminiproblems-problemphyxmini0834-usesvacuumcoulombconstant, def:physics:phyx-mini-0834:phyxminiproblems-problemphyxmini0834-hasphysicalrectangleparameters, def:physics:phyx-mini-0834:phyxminiproblems-problemphyxmini0834-satisfiespointchargecoulombfieldlaw, def:physics:phyx-mini-0834:phyxminiproblems-problemphyxmini0834-satisfieselectricfieldsuperposition, def:physics:phyx-mini-0834:phyxminiproblems-problemphyxmini0834-resultantmagnituderepresentsfieldatdot, def:physics:phyx-mini-0834:phyxminiproblems-problemphyxmini0834-recordeddatasetanswer, def:physics:phyx-mini-0834:phyxminiproblems-problemphyxmini0834-roundstonearestthousand, def:physics:phyx-mini-0834:phyxminiproblems-problemphyxmini0834-isuniquenearestdisplayedfieldmagnitude}
The assigned autoformalize agent should translate this physics problem into Lean declarations in the covered file, with theorem and lemma proof bodies written as `by sorry`.
\end{theorem}
\begin{proof}
This is an autoformalization task, not a proof task. Produce statements that can later be proved by the physics prover without weakening the physical model.
\end{proof}

% --- Archon declaration topology begin ---
\section{Declaration topology}
\begin{definition}[Length Quantity]
  \label{def:physics:phyx-mini-0834:phyxminiproblems-problemphyxmini0834-lengthquantity}
  \lean{PhyXMiniProblems.ProblemPhyXMini0834.LengthQuantity}
  A nonnegative, unit-independent physical length
\end{definition}
\begin{proof}
  This is the declared model definition of Length Quantity.
\end{proof}
\begin{definition}[Signed Charge Quantity]
  \label{def:physics:phyx-mini-0834:phyxminiproblems-problemphyxmini0834-signedchargequantity}
  \lean{PhyXMiniProblems.ProblemPhyXMini0834.SignedChargeQuantity}
  A signed, unit-independent physical electric charge
\end{definition}
\begin{proof}
  This is the declared model definition of Signed Charge Quantity.
\end{proof}
\begin{definition}[electric Field Strength Dimension]
  \label{def:physics:phyx-mini-0834:phyxminiproblems-problemphyxmini0834-electricfieldstrengthdimension}
  \lean{PhyXMiniProblems.ProblemPhyXMini0834.electricFieldStrengthDimension}
  The physical dimension M L T⁻² C⁻¹ of electric-field strength
\end{definition}
\begin{proof}
  This is the declared model definition of electric Field Strength Dimension.
\end{proof}
\begin{definition}[Electric Field Magnitude Quantity]
  \label{def:physics:phyx-mini-0834:phyxminiproblems-problemphyxmini0834-electricfieldmagnitudequantity}
  \lean{PhyXMiniProblems.ProblemPhyXMini0834.ElectricFieldMagnitudeQuantity}
  \uses{def:physics:phyx-mini-0834:phyxminiproblems-problemphyxmini0834-electricfieldstrengthdimension}
  A nonnegative physical electric-field magnitude
\end{definition}
\begin{proof}
  This is the declared model definition of Electric Field Magnitude Quantity.
\end{proof}
\begin{definition}[length In Meters]
  \label{def:physics:phyx-mini-0834:phyxminiproblems-problemphyxmini0834-lengthinmeters}
  \lean{PhyXMiniProblems.ProblemPhyXMini0834.lengthInMeters}
  \uses{def:physics:phyx-mini-0834:phyxminiproblems-problemphyxmini0834-lengthquantity}
  Coherent-SI metre readout of a physical length
\end{definition}
\begin{proof}
  This is the declared model definition of length In Meters.
\end{proof}
\begin{definition}[length In Centimeters]
  \label{def:physics:phyx-mini-0834:phyxminiproblems-problemphyxmini0834-lengthincentimeters}
  \lean{PhyXMiniProblems.ProblemPhyXMini0834.lengthInCentimeters}
  \uses{def:physics:phyx-mini-0834:phyxminiproblems-problemphyxmini0834-lengthquantity, def:physics:phyx-mini-0834:phyxminiproblems-problemphyxmini0834-lengthinmeters}
  Centimetre readout used by the two dimension labels in the image
\end{definition}
\begin{proof}
  This is the declared model definition of length In Centimeters.
\end{proof}
\begin{definition}[charge In Coulombs]
  \label{def:physics:phyx-mini-0834:phyxminiproblems-problemphyxmini0834-chargeincoulombs}
  \lean{PhyXMiniProblems.ProblemPhyXMini0834.chargeInCoulombs}
  \uses{def:physics:phyx-mini-0834:phyxminiproblems-problemphyxmini0834-signedchargequantity}
  Coherent-SI coulomb readout of a signed physical charge
\end{definition}
\begin{proof}
  This is the declared model definition of charge In Coulombs.
\end{proof}
\begin{definition}[charge In Nanocoulombs]
  \label{def:physics:phyx-mini-0834:phyxminiproblems-problemphyxmini0834-chargeinnanocoulombs}
  \lean{PhyXMiniProblems.ProblemPhyXMini0834.chargeInNanocoulombs}
  \uses{def:physics:phyx-mini-0834:phyxminiproblems-problemphyxmini0834-signedchargequantity, def:physics:phyx-mini-0834:phyxminiproblems-problemphyxmini0834-chargeincoulombs}
  Nanocoulomb readout used by the three printed charge labels
\end{definition}
\begin{proof}
  This is the declared model definition of charge In Nanocoulombs.
\end{proof}
\begin{definition}[field Magnitude In Newtons Per Coulomb]
  \label{def:physics:phyx-mini-0834:phyxminiproblems-problemphyxmini0834-fieldmagnitudeinnewtonspercoulomb}
  \lean{PhyXMiniProblems.ProblemPhyXMini0834.fieldMagnitudeInNewtonsPerCoulomb}
  \uses{def:physics:phyx-mini-0834:phyxminiproblems-problemphyxmini0834-electricfieldmagnitudequantity}
  Coherent-SI readout of a field magnitude in newtons per coulomb
\end{definition}
\begin{proof}
  This is the declared model definition of field Magnitude In Newtons Per Coulomb.
\end{proof}
\begin{definition}[Rectangle Corner]
  \label{def:physics:phyx-mini-0834:phyxminiproblems-problemphyxmini0834-rectanglecorner}
  \lean{PhyXMiniProblems.ProblemPhyXMini0834.RectangleCorner}
  The four corners of the dashed rectangle in the supplied figure
\end{definition}
\begin{proof}
  This is the declared model definition of Rectangle Corner.
\end{proof}
\begin{definition}[Source Charge]
  \label{def:physics:phyx-mini-0834:phyxminiproblems-problemphyxmini0834-sourcecharge}
  \lean{PhyXMiniProblems.ProblemPhyXMini0834.SourceCharge}
  The three point charges, distinguished by their image locations
\end{definition}
\begin{proof}
  This is the declared model definition of Source Charge.
\end{proof}
\begin{definition}[expected Source Corner]
  \label{def:physics:phyx-mini-0834:phyxminiproblems-problemphyxmini0834-expectedsourcecorner}
  \lean{PhyXMiniProblems.ProblemPhyXMini0834.expectedSourceCorner}
  \uses{def:physics:phyx-mini-0834:phyxminiproblems-problemphyxmini0834-rectanglecorner, def:physics:phyx-mini-0834:phyxminiproblems-problemphyxmini0834-sourcecharge}
  The corner at which each named source appears in the primary image
\end{definition}
\begin{proof}
  This is the declared model definition of expected Source Corner.
\end{proof}
\begin{definition}[Figure Charge Color]
  \label{def:physics:phyx-mini-0834:phyxminiproblems-problemphyxmini0834-figurechargecolor}
  \lean{PhyXMiniProblems.ProblemPhyXMini0834.FigureChargeColor}
  The teal and red fill colours used to distinguish charge signs
\end{definition}
\begin{proof}
  This is the declared model definition of Figure Charge Color.
\end{proof}
\begin{definition}[Figure Charge Sign]
  \label{def:physics:phyx-mini-0834:phyxminiproblems-problemphyxmini0834-figurechargesign}
  \lean{PhyXMiniProblems.ProblemPhyXMini0834.FigureChargeSign}
  The sign glyph drawn inside a charge circle
\end{definition}
\begin{proof}
  This is the declared model definition of Figure Charge Sign.
\end{proof}
\begin{definition}[expected Source Color]
  \label{def:physics:phyx-mini-0834:phyxminiproblems-problemphyxmini0834-expectedsourcecolor}
  \lean{PhyXMiniProblems.ProblemPhyXMini0834.expectedSourceColor}
  \uses{def:physics:phyx-mini-0834:phyxminiproblems-problemphyxmini0834-sourcecharge, def:physics:phyx-mini-0834:phyxminiproblems-problemphyxmini0834-figurechargecolor}
  Colour of each source circle in the supplied raster
\end{definition}
\begin{proof}
  This is the declared model definition of expected Source Color.
\end{proof}
\begin{definition}[expected Source Sign]
  \label{def:physics:phyx-mini-0834:phyxminiproblems-problemphyxmini0834-expectedsourcesign}
  \lean{PhyXMiniProblems.ProblemPhyXMini0834.expectedSourceSign}
  \uses{def:physics:phyx-mini-0834:phyxminiproblems-problemphyxmini0834-sourcecharge, def:physics:phyx-mini-0834:phyxminiproblems-problemphyxmini0834-figurechargesign}
  Sign glyph of each source circle in the supplied raster
\end{definition}
\begin{proof}
  This is the declared model definition of expected Source Sign.
\end{proof}
\begin{definition}[expected Charge In Nanocoulombs]
  \label{def:physics:phyx-mini-0834:phyxminiproblems-problemphyxmini0834-expectedchargeinnanocoulombs}
  \lean{PhyXMiniProblems.ProblemPhyXMini0834.expectedChargeInNanocoulombs}
  \uses{def:physics:phyx-mini-0834:phyxminiproblems-problemphyxmini0834-sourcecharge}
  Signed nanocoulomb value printed beside each source
\end{definition}
\begin{proof}
  This is the declared model definition of expected Charge In Nanocoulombs.
\end{proof}
\begin{definition}[Figure Dimension Side]
  \label{def:physics:phyx-mini-0834:phyxminiproblems-problemphyxmini0834-figuredimensionside}
  \lean{PhyXMiniProblems.ProblemPhyXMini0834.FigureDimensionSide}
  The two sides carrying numerical dimension labels
\end{definition}
\begin{proof}
  This is the declared model definition of Figure Dimension Side.
\end{proof}
\begin{definition}[Rectangular Charge Figure]
  \label{def:physics:phyx-mini-0834:phyxminiproblems-problemphyxmini0834-rectangularchargefigure}
  \lean{PhyXMiniProblems.ProblemPhyXMini0834.RectangularChargeFigure}
  \uses{def:physics:phyx-mini-0834:phyxminiproblems-problemphyxmini0834-rectanglecorner, def:physics:phyx-mini-0834:phyxminiproblems-problemphyxmini0834-sourcecharge, def:physics:phyx-mini-0834:phyxminiproblems-problemphyxmini0834-figurechargecolor, def:physics:phyx-mini-0834:phyxminiproblems-problemphyxmini0834-figurechargesign, def:physics:phyx-mini-0834:phyxminiproblems-problemphyxmini0834-figuredimensionside}
  Literal presentation data transcribed from image 834.png
\end{definition}
\begin{proof}
  This is the declared model definition of Rectangular Charge Figure.
\end{proof}
\begin{definition}[position Vector In Meters]
  \label{def:physics:phyx-mini-0834:phyxminiproblems-problemphyxmini0834-positionvectorinmeters}
  \lean{PhyXMiniProblems.ProblemPhyXMini0834.positionVectorInMeters}
  Turn a planar Space 2 point into its explicitly metre-valued vector
\end{definition}
\begin{proof}
  This is the declared model definition of position Vector In Meters.
\end{proof}
\begin{definition}[Rectangular Point Charge Setup]
  \label{def:physics:phyx-mini-0834:phyxminiproblems-problemphyxmini0834-rectangularpointchargesetup}
  \lean{PhyXMiniProblems.ProblemPhyXMini0834.RectangularPointChargeSetup}
  \uses{def:physics:phyx-mini-0834:phyxminiproblems-problemphyxmini0834-lengthquantity, def:physics:phyx-mini-0834:phyxminiproblems-problemphyxmini0834-signedchargequantity, def:physics:phyx-mini-0834:phyxminiproblems-problemphyxmini0834-electricfieldmagnitudequantity, def:physics:phyx-mini-0834:phyxminiproblems-problemphyxmini0834-rectanglecorner, def:physics:phyx-mini-0834:phyxminiproblems-problemphyxmini0834-sourcecharge, def:physics:phyx-mini-0834:phyxminiproblems-problemphyxmini0834-rectangularchargefigure}
  Independent physical quantities and fields in the electrostatic setup. The resultant magnitude is an unknown physical observable. It is not defined from a displayed answer; a separate premise below relates it to the norm of the independently modeled resultant vector field
\end{definition}
\begin{proof}
  This is the declared model definition of Rectangular Point Charge Setup.
\end{proof}
\begin{definition}[source Position]
  \label{def:physics:phyx-mini-0834:phyxminiproblems-problemphyxmini0834-sourceposition}
  \lean{PhyXMiniProblems.ProblemPhyXMini0834.sourcePosition}
  \uses{def:physics:phyx-mini-0834:phyxminiproblems-problemphyxmini0834-sourcecharge, def:physics:phyx-mini-0834:phyxminiproblems-problemphyxmini0834-rectangularpointchargesetup}
  Spatial position of a named source charge
\end{definition}
\begin{proof}
  This is the declared model definition of source Position.
\end{proof}
\begin{definition}[observation Position]
  \label{def:physics:phyx-mini-0834:phyxminiproblems-problemphyxmini0834-observationposition}
  \lean{PhyXMiniProblems.ProblemPhyXMini0834.observationPosition}
  \uses{def:physics:phyx-mini-0834:phyxminiproblems-problemphyxmini0834-rectangularpointchargesetup}
  Spatial position marked by the black observation dot
\end{definition}
\begin{proof}
  This is the declared model definition of observation Position.
\end{proof}
\begin{definition}[displacement From Source In Meters]
  \label{def:physics:phyx-mini-0834:phyxminiproblems-problemphyxmini0834-displacementfromsourceinmeters}
  \lean{PhyXMiniProblems.ProblemPhyXMini0834.displacementFromSourceInMeters}
  \uses{def:physics:phyx-mini-0834:phyxminiproblems-problemphyxmini0834-sourcecharge, def:physics:phyx-mini-0834:phyxminiproblems-problemphyxmini0834-positionvectorinmeters, def:physics:phyx-mini-0834:phyxminiproblems-problemphyxmini0834-rectangularpointchargesetup, def:physics:phyx-mini-0834:phyxminiproblems-problemphyxmini0834-sourceposition}
  Displacement from a named charge to an arbitrary field point, in metres
\end{definition}
\begin{proof}
  This is the declared model definition of displacement From Source In Meters.
\end{proof}
\begin{definition}[resultant Field At Dot In Newtons Per Coulomb]
  \label{def:physics:phyx-mini-0834:phyxminiproblems-problemphyxmini0834-resultantfieldatdotinnewtonspercoulomb}
  \lean{PhyXMiniProblems.ProblemPhyXMini0834.resultantFieldAtDotInNewtonsPerCoulomb}
  \uses{def:physics:phyx-mini-0834:phyxminiproblems-problemphyxmini0834-rectangularpointchargesetup, def:physics:phyx-mini-0834:phyxminiproblems-problemphyxmini0834-observationposition}
  Resultant field vector at the black dot, in newtons per coulomb
\end{definition}
\begin{proof}
  This is the declared model definition of resultant Field At Dot In Newtons Per Coulomb.
\end{proof}
\begin{definition}[resultant Field Norm At Dot In Newtons Per Coulomb]
  \label{def:physics:phyx-mini-0834:phyxminiproblems-problemphyxmini0834-resultantfieldnormatdotinnewtonspercoulomb}
  \lean{PhyXMiniProblems.ProblemPhyXMini0834.resultantFieldNormAtDotInNewtonsPerCoulomb}
  \uses{def:physics:phyx-mini-0834:phyxminiproblems-problemphyxmini0834-rectangularpointchargesetup, def:physics:phyx-mini-0834:phyxminiproblems-problemphyxmini0834-resultantfieldatdotinnewtonspercoulomb}
  Norm of the resultant vector field at the dot, in newtons per coulomb
\end{definition}
\begin{proof}
  This is the declared model definition of resultant Field Norm At Dot In Newtons Per Coulomb.
\end{proof}
\begin{definition}[Matches Supplied Charge Rectangle Figure]
  \label{def:physics:phyx-mini-0834:phyxminiproblems-problemphyxmini0834-matchessuppliedchargerectanglefigure}
  \lean{PhyXMiniProblems.ProblemPhyXMini0834.MatchesSuppliedChargeRectangleFigure}
  \uses{def:physics:phyx-mini-0834:phyxminiproblems-problemphyxmini0834-lengthinmeters, def:physics:phyx-mini-0834:phyxminiproblems-problemphyxmini0834-lengthincentimeters, def:physics:phyx-mini-0834:phyxminiproblems-problemphyxmini0834-chargeinnanocoulombs, def:physics:phyx-mini-0834:phyxminiproblems-problemphyxmini0834-expectedsourcecorner, def:physics:phyx-mini-0834:phyxminiproblems-problemphyxmini0834-expectedsourcecolor, def:physics:phyx-mini-0834:phyxminiproblems-problemphyxmini0834-expectedsourcesign, def:physics:phyx-mini-0834:phyxminiproblems-problemphyxmini0834-expectedchargeinnanocoulombs, def:physics:phyx-mini-0834:phyxminiproblems-problemphyxmini0834-positionvectorinmeters, def:physics:phyx-mini-0834:phyxminiproblems-problemphyxmini0834-rectangularpointchargesetup}
  Literal image evidence, numerical labels, and their association with the physical charges and rectangle dimensions. No resultant field value appears here
\end{definition}
\begin{proof}
  This is the declared model definition of Matches Supplied Charge Rectangle Figure.
\end{proof}
\begin{definition}[Uses Vacuum Coulomb Constant]
  \label{def:physics:phyx-mini-0834:phyxminiproblems-problemphyxmini0834-usesvacuumcoulombconstant}
  \lean{PhyXMiniProblems.ProblemPhyXMini0834.UsesVacuumCoulombConstant}
  \uses{def:physics:phyx-mini-0834:phyxminiproblems-problemphyxmini0834-rectangularpointchargesetup}
  Independent free-space calibration of Physlib's electromagnetic system. The scalar has coherent-SI units N m²/C² in the Coulomb formula below
\end{definition}
\begin{proof}
  This is the declared model definition of Uses Vacuum Coulomb Constant.
\end{proof}
\begin{definition}[Has Physical Rectangle Parameters]
  \label{def:physics:phyx-mini-0834:phyxminiproblems-problemphyxmini0834-hasphysicalrectangleparameters}
  \lean{PhyXMiniProblems.ProblemPhyXMini0834.HasPhysicalRectangleParameters}
  \uses{def:physics:phyx-mini-0834:phyxminiproblems-problemphyxmini0834-lengthinmeters, def:physics:phyx-mini-0834:phyxminiproblems-problemphyxmini0834-rectangularpointchargesetup, def:physics:phyx-mini-0834:phyxminiproblems-problemphyxmini0834-observationposition, def:physics:phyx-mini-0834:phyxminiproblems-problemphyxmini0834-displacementfromsourceinmeters}
  Positivity and separation conditions selecting the physical branch
\end{definition}
\begin{proof}
  This is the declared model definition of Has Physical Rectangle Parameters.
\end{proof}
\begin{definition}[Satisfies Point Charge Coulomb Field Law]
  \label{def:physics:phyx-mini-0834:phyxminiproblems-problemphyxmini0834-satisfiespointchargecoulombfieldlaw}
  \lean{PhyXMiniProblems.ProblemPhyXMini0834.SatisfiesPointChargeCoulombFieldLaw}
  \uses{def:physics:phyx-mini-0834:phyxminiproblems-problemphyxmini0834-chargeincoulombs, def:physics:phyx-mini-0834:phyxminiproblems-problemphyxmini0834-rectangularpointchargesetup, def:physics:phyx-mini-0834:phyxminiproblems-problemphyxmini0834-sourceposition, def:physics:phyx-mini-0834:phyxminiproblems-problemphyxmini0834-displacementfromsourceinmeters}
  For a source charge q and displacement vector r from the source to the field point, the electric field is k q r / ‖r‖³. This all-points law is stated only away from the point source and contains no requested result
\end{definition}
\begin{proof}
  This is the declared model definition of Satisfies Point Charge Coulomb Field Law.
\end{proof}
\begin{definition}[Satisfies Electric Field Superposition]
  \label{def:physics:phyx-mini-0834:phyxminiproblems-problemphyxmini0834-satisfieselectricfieldsuperposition}
  \lean{PhyXMiniProblems.ProblemPhyXMini0834.SatisfiesElectricFieldSuperposition}
  \uses{def:physics:phyx-mini-0834:phyxminiproblems-problemphyxmini0834-sourcecharge, def:physics:phyx-mini-0834:phyxminiproblems-problemphyxmini0834-rectangularpointchargesetup}
  The total electric field is the vector sum of all three source fields
\end{definition}
\begin{proof}
  This is the declared model definition of Satisfies Electric Field Superposition.
\end{proof}
\begin{definition}[Resultant Magnitude Represents Field At Dot]
  \label{def:physics:phyx-mini-0834:phyxminiproblems-problemphyxmini0834-resultantmagnituderepresentsfieldatdot}
  \lean{PhyXMiniProblems.ProblemPhyXMini0834.ResultantMagnitudeRepresentsFieldAtDot}
  \uses{def:physics:phyx-mini-0834:phyxminiproblems-problemphyxmini0834-fieldmagnitudeinnewtonspercoulomb, def:physics:phyx-mini-0834:phyxminiproblems-problemphyxmini0834-rectangularpointchargesetup, def:physics:phyx-mini-0834:phyxminiproblems-problemphyxmini0834-resultantfieldnormatdotinnewtonspercoulomb}
  The dimensionful scalar observable is the Euclidean magnitude of the modeled resultant vector at the black dot. This states a measurement relation, not a numerical answer
\end{definition}
\begin{proof}
  This is the declared model definition of Resultant Magnitude Represents Field At Dot.
\end{proof}
\begin{definition}[Answer Choice]
  \label{def:physics:phyx-mini-0834:phyxminiproblems-problemphyxmini0834-answerchoice}
  \lean{PhyXMiniProblems.ProblemPhyXMini0834.AnswerChoice}
  Labels attached to the four multiple-choice answers
\end{definition}
\begin{proof}
  This is the declared model definition of Answer Choice.
\end{proof}
\begin{definition}[field Magnitude In Newtons Per Coulomb]
  \label{def:physics:phyx-mini-0834:phyxminiproblems-problemphyxmini0834-answerchoice-fieldmagnitudeinnewtonspercoulomb}
  \lean{PhyXMiniProblems.ProblemPhyXMini0834.AnswerChoice.fieldMagnitudeInNewtonsPerCoulomb}
  \uses{def:physics:phyx-mini-0834:phyxminiproblems-problemphyxmini0834-fieldmagnitudeinnewtonspercoulomb, def:physics:phyx-mini-0834:phyxminiproblems-problemphyxmini0834-answerchoice}
  Electric-field magnitude in newtons per coulomb printed by each choice
\end{definition}
\begin{proof}
  This is the declared model definition of field Magnitude In Newtons Per Coulomb.
\end{proof}
\begin{definition}[recorded Dataset Answer]
  \label{def:physics:phyx-mini-0834:phyxminiproblems-problemphyxmini0834-recordeddatasetanswer}
  \lean{PhyXMiniProblems.ProblemPhyXMini0834.recordedDatasetAnswer}
  \uses{def:physics:phyx-mini-0834:phyxminiproblems-problemphyxmini0834-answerchoice}
  The answer label recorded by the supplied dataset metadata
\end{definition}
\begin{proof}
  This is the declared model definition of recorded Dataset Answer.
\end{proof}
\begin{definition}[Rounds To Nearest Thousand]
  \label{def:physics:phyx-mini-0834:phyxminiproblems-problemphyxmini0834-roundstonearestthousand}
  \lean{PhyXMiniProblems.ProblemPhyXMini0834.RoundsToNearestThousand}
  A value rounds to a displayed reading at the nearest 10³ N/C
\end{definition}
\begin{proof}
  This is the declared model definition of Rounds To Nearest Thousand.
\end{proof}
\begin{definition}[Is Unique Nearest Displayed Field Magnitude]
  \label{def:physics:phyx-mini-0834:phyxminiproblems-problemphyxmini0834-isuniquenearestdisplayedfieldmagnitude}
  \lean{PhyXMiniProblems.ProblemPhyXMini0834.IsUniqueNearestDisplayedFieldMagnitude}
  \uses{def:physics:phyx-mini-0834:phyxminiproblems-problemphyxmini0834-fieldmagnitudeinnewtonspercoulomb, def:physics:phyx-mini-0834:phyxminiproblems-problemphyxmini0834-answerchoice}
  A displayed choice is strictly nearer than every alternative
\end{definition}
\begin{proof}
  This is the declared model definition of Is Unique Nearest Displayed Field Magnitude.
\end{proof}
\begin{lemma}[resultant Field At Dot eq coulomb Sum]
  \label{lem:physics:phyx-mini-0834:phyxminiproblems-problemphyxmini0834-resultantfieldatdot-eq-coulombsum}
  \lean{PhyXMiniProblems.ProblemPhyXMini0834.resultantFieldAtDot_eq_coulombSum}
  \uses{def:physics:phyx-mini-0834:phyxminiproblems-problemphyxmini0834-chargeincoulombs, def:physics:phyx-mini-0834:phyxminiproblems-problemphyxmini0834-sourcecharge, def:physics:phyx-mini-0834:phyxminiproblems-problemphyxmini0834-rectangularpointchargesetup, def:physics:phyx-mini-0834:phyxminiproblems-problemphyxmini0834-observationposition, def:physics:phyx-mini-0834:phyxminiproblems-problemphyxmini0834-displacementfromsourceinmeters, def:physics:phyx-mini-0834:phyxminiproblems-problemphyxmini0834-resultantfieldatdotinnewtonspercoulomb, def:physics:phyx-mini-0834:phyxminiproblems-problemphyxmini0834-hasphysicalrectangleparameters, def:physics:phyx-mini-0834:phyxminiproblems-problemphyxmini0834-satisfiespointchargecoulombfieldlaw, def:physics:phyx-mini-0834:phyxminiproblems-problemphyxmini0834-satisfieselectricfieldsuperposition}
  Coulomb's vector law and superposition give the source-by-source expression for the field at the observation dot. The expression contains only the independent source data and governing constant
\end{lemma}
\begin{proof}
  The conclusion follows from the governing relations and figure readouts named in the statement of resultant Field At Dot eq coulomb Sum.
\end{proof}
\begin{lemma}[resultant Field Norm At Dot numerical Bounds]
  \label{lem:physics:phyx-mini-0834:phyxminiproblems-problemphyxmini0834-resultantfieldnormatdot-numericalbounds}
  \lean{PhyXMiniProblems.ProblemPhyXMini0834.resultantFieldNormAtDot_numericalBounds}
  \uses{def:physics:phyx-mini-0834:phyxminiproblems-problemphyxmini0834-rectangularpointchargesetup, def:physics:phyx-mini-0834:phyxminiproblems-problemphyxmini0834-resultantfieldnormatdotinnewtonspercoulomb, def:physics:phyx-mini-0834:phyxminiproblems-problemphyxmini0834-matchessuppliedchargerectanglefigure, def:physics:phyx-mini-0834:phyxminiproblems-problemphyxmini0834-usesvacuumcoulombconstant, def:physics:phyx-mini-0834:phyxminiproblems-problemphyxmini0834-hasphysicalrectangleparameters, def:physics:phyx-mini-0834:phyxminiproblems-problemphyxmini0834-satisfiespointchargecoulombfieldlaw, def:physics:phyx-mini-0834:phyxminiproblems-problemphyxmini0834-satisfieselectricfieldsuperposition}
  With the three image charges and rectangle dimensions, the resultant magnitude lies between 85.5 kN/C and 86.5 kN/C. In particular it rounds to the displayed 8.6 * 10\textasciicircum{}4 N/C value
\end{lemma}
\begin{proof}
  The conclusion follows from the governing relations and figure readouts named in the statement of resultant Field Norm At Dot numerical Bounds.
\end{proof}
% --- Archon declaration topology end ---
% --- Archon physics formalization source end ---
